\section{用微分与 Frobenius 内积读取梯度}
\label{sec:10-Matrix-Calculus/02-Differentials/01}

手头的任务是最小化

\[
f(X)=\frac12\norm{AX-B}_F^2,
\qquad
A\in\R^{100\times 50},\quad X\in\R^{50\times 30},\quad B\in\R^{100\times 30}.
\]

梯度下降要的是 \(\nabla_X f\)，一个 \(50\times 30\) 的矩阵，里面有 \(1500\) 个偏导数。没有人会一个个去求它们再拼回形状。

一条看起来正规的出路是把矩阵拉平。令 \(x=\operatorname{vec}(X)\in\R^{1500}\)，用 Kronecker 积把问题写成普通的向量最小二乘，梯度是 \((I_{30}\otimes A^{\trans}A)x-(I_{30}\otimes A^{\trans})\operatorname{vec}(B)\)。这个式子没错，可它把一个 \(1500\times 1500\) 的对象摆在了你面前，而 \(A\)、\(X\)、\(B\) 各自的形状全被拉平抹掉了。你再也没法从结果里一眼认出 \(A^{\trans}(AX-B)\)，也没法判断某一项该不该转置。真正在代码里要写的那一行，反而看不出来了。

需要的是一套既保住矩阵形状、又能机械推演的记法。微分符号 \(\mathrm dX\) 加上 Frobenius 内积就是干这个的。它把``求梯度''换成一件纯代数的活：把 \(\mathrm df\) 整理成 \(\ip{G}{\mathrm dX}_F\) 的样子，括号里的 \(G\) 就是梯度。

\begin{center}
\begin{tikzpicture}[x=1cm,y=1cm,font=\footnotesize,>=stealth,
    box/.style={draw=black!55,rounded corners=2pt,inner sep=4pt,align=center}]
  \node[box,fill=blue!6] (f)  at (0.9,1.35) {\(f(X)\)};
  \node[box,fill=blue!6] (t1) at (3.7,2.70) {写出微分\\\(\mathrm df\) 对 \(\mathrm dX\) 线性};
  \node[box,fill=blue!6] (t2) at (7.2,2.70) {整理成内积\\\(\mathrm df=\ip{G}{\mathrm dX}_F\)};
  \node[box,fill=blue!6] (t3) at (10.2,2.70) {读出\\\(\nabla_X f=G\)};
  \node[box,fill=black!4] (b1) at (3.7,0.00) {逐元素求\\\(\partial f/\partial X_{ij}\)};
  \node[box,fill=black!4] (b2) at (7.2,0.00) {\(1500\) 个标量\\与一堆求和指标};
  \node[box,fill=black!4] (b3) at (10.2,0.00) {按 \((i,j)\) 拼回形状};
  \draw[->,blue!55!black,thick] (f) -- (t1);
  \draw[->,blue!55!black,thick] (t1) -- (t2);
  \draw[->,blue!55!black,thick] (t2) -- (t3);
  \draw[->,black!45,dashed] (f) -- (b1);
  \draw[->,black!45,dashed] (b1) -- (b2);
  \draw[->,black!45,dashed] (b2) -- (b3);
\end{tikzpicture}

\emph{上下两条路通向同一个矩阵；上面一条从头到尾没有出现过下标。}
\end{center}

图上那条实线是本节要走的路线，三步都是可以机械执行的：套微分规则、用迹的循环性质把 \(\mathrm dX\) 挪到最右端、把它前面那一坨转置一下读出来。下面那条虚线不是错的，只是每一步都在跟指标搏斗，而且中途丢掉的形状信息到最后还得靠人脑补回去。下面把实线的三步逐个说清楚。

\subsection{微分是线性化，\(\mathrm dX\) 只是占位符}

先把 \(\mathrm dX\) 的身份说定。它不是``一个极小的矩阵''，也不需要任何关于无穷小的哲学。它是一个与 \(X\) 同形的符号，代表``扰动发生在哪个位置、朝哪个方向''，而 \(\mathrm df\) 是 \(f\) 相应变化的线性主部：

\[
f(X+\mathrm dX)-f(X)=\mathrm df+o(\norm{\mathrm dX}_F),
\]

其中 \(\mathrm df\) 对 \(\mathrm dX\) 线性。这与上一单元的 \(Df(X)[H]\) 是同一件事，换了个书写习惯：把 \(H\) 写成 \(\mathrm dX\)，好处是可以在一行式子里连着用乘法、转置、迹的规则往下推。

规则只有三条，都从代入展开、扔掉二阶项得来：

\[
\mathrm d(A+B)=\mathrm dA+\mathrm dB,
\qquad
\mathrm d(AB)=(\mathrm dA)B+A\,\mathrm dB,
\qquad
\mathrm d(A^{\trans})=(\mathrm dA)^{\trans}.
\]

乘积那条要留意顺序不能动，矩阵乘法不可交换。取 \(Y=X^{\trans}X\) 试一次：代入 \(X+\mathrm dX\) 展开得 \(X^{\trans}X+(\mathrm dX)^{\trans}X+X^{\trans}\mathrm dX+(\mathrm dX)^{\trans}\mathrm dX\)，末项是二阶量，去掉后

\[
\mathrm dY=(\mathrm dX)^{\trans}X+X^{\trans}\mathrm dX.
\]

条件对一遍。``扔掉二阶项''之所以合法，是因为 Fréchet 可微性要求余项除以 \(\norm{\mathrm dX}_F\) 后趋于零，而 \((\mathrm dX)^{\trans}\mathrm dX\) 的 Frobenius 范数不超过 \(\norm{\mathrm dX}_F^2\)，正好满足。对于由加、乘、转置、迹、逆、行列式复合出来的函数，这一步永远成立；碰到 \(\max\)、绝对值、取整这类地方就要停下来单独看，那里的可微性本身是有条件的。

\subsection{Frobenius 内积把梯度钉成唯一的那个矩阵}

两个同形矩阵的 Frobenius 内积是逐元素乘积求和：

\[
\ip{A}{B}_F=\sum_{i,j}A_{ij}B_{ij}=\trace(A^{\trans}B).
\]

它对称、双线性、正定，矩阵空间在它之下就是一个普通的欧氏空间，``两个矩阵的距离''``投影到某个子空间''这类说法才有依据。

梯度的定义随之而来：\(\nabla_X f\) 是使

\[
\mathrm df=\ip{\nabla_X f}{\mathrm dX}_F=\trace\bigl((\nabla_X f)^{\trans}\mathrm dX\bigr)
\]

对一切 \(\mathrm dX\) 成立的那个矩阵。内积的正定性保证这样的矩阵至多一个：若两个矩阵都满足，它们的差与所有 \(\mathrm dX\) 内积为零，取 \(\mathrm dX\) 等于这个差本身即得差为零。唯一性是这套方法能用的前提——只要你把 \(\mathrm df\) 整理成了内积形式，前面那个矩阵就是梯度，没有第二种可能，也不必回头验算。

\subsection{一个完整范例走完三步}

回到开头的最小二乘。记残差 \(R=AX-B\)，于是 \(f=\tfrac12\trace(R^{\trans}R)\)。

第一步，写微分。乘积规则给出 \(\mathrm df=\tfrac12\trace\bigl((\mathrm dR)^{\trans}R+R^{\trans}\mathrm dR\bigr)\)。迹在转置下不变，两项相等，因子 \(\tfrac12\) 被吃掉：

\[
\mathrm df=\trace(R^{\trans}\mathrm dR).
\]

顺带解释一下那个 \(\tfrac12\) 是怎么来的。它是建模者自己加的归一化，不加的话微分里会留下一个 \(2\)，梯度表达式多一个因子；最优解不受影响。这类常数没有定理可背，只有约定可查。

第二步，把 \(\mathrm dR\) 换成 \(\mathrm dX\)。\(A\) 与 \(B\) 是常量，\(\mathrm dA=0\)、\(\mathrm dB=0\)，故 \(\mathrm dR=A\,\mathrm dX\)，代入得 \(\mathrm df=\trace(R^{\trans}A\,\mathrm dX)\)。

第三步，读取。为了对上 \(\trace(G^{\trans}\mathrm dX)\)，把 \(R^{\trans}A\) 写成某个矩阵的转置：\(R^{\trans}A=(A^{\trans}R)^{\trans}\)，于是

\[
\mathrm df=\trace\bigl((A^{\trans}R)^{\trans}\mathrm dX\bigr),
\qquad
\nabla_X f=A^{\trans}(AX-B).
\]

形状核对：\(A^{\trans}\) 是 \(50\times 100\)，\(AX-B\) 是 \(100\times 30\)，乘出来 \(50\times 30\)，与 \(X\) 同形。跟前面那个 Kronecker 版本对照一下，同样的信息，这边可以直接抄进代码，而且每个因子为什么在那个位置都看得见。

\subsection{最后一步的转置是最容易漏的地方}

能挪动因子的操作只有迹的循环性质：

\[
\trace(ABC)=\trace(BCA)=\trace(CAB).
\]

它是整体轮换，不是任意交换，\(\trace(ABC)\) 一般不等于 \(\trace(ACB)\)。轮换配合转置，可以把内积里的因子从一侧搬到另一侧：

\[
\ip{A}{BC}_F=\trace(A^{\trans}BC)=\trace(CA^{\trans}B)
=\trace\bigl((AC^{\trans})^{\trans}B\bigr)=\ip{AC^{\trans}}{B}_F.
\]

推导中反复用到的就是这个动作：把 \(\mathrm dX\) 从表达式中间挪到最右端，路上遇到的因子被吞进梯度的转置里面。

出错最集中的是最后那一下。假设整理到了

\[
\mathrm df=\trace(BA\,\mathrm dX),
\]

与 \(\trace(G^{\trans}\mathrm dX)\) 比较，得到的是 \(G^{\trans}=BA\)，因此 \(G=A^{\trans}B^{\trans}\)，而不是 \(BA\)。把 \(\mathrm dX\) 前面那一坨直接叫作梯度，是这套方法里最高频的失误，尤其在 \(BA\) 恰好是方阵时形状检查还抓不住。可靠的做法是每次都把定义式 \(\trace\bigl((\nabla_X f)^{\trans}\mathrm dX\bigr)\) 完整写出来再比对，不靠视觉记忆决定要不要转置。

\subsection{推导一个新层的梯度不必展开到元素}

这套流程真正省力的场合是给新结构写反向传播。设某一层的输出是 \(Y=\sigma(XW)\)，下游传来的量是 \(\nabla_Y L\)，你要的是 \(\nabla_W L\) 和 \(\nabla_X L\)。按流程走：先写 \(\mathrm dL=\ip{\nabla_Y L}{\mathrm dY}_F\)，再把 \(\mathrm dY=\sigma'(XW)\odot(\mathrm dX\,W+X\,\mathrm dW)\) 代进去。逐元素乘 \(\odot\) 与 Frobenius 内积之间有一条现成的搬移规则 \(\ip{A}{B\odot C}_F=\ip{A\odot B}{C}_F\)，用它把 \(\sigma'\) 挪到左边，记 \(\Delta=\nabla_Y L\odot\sigma'(XW)\)，剩下的就是

\[
\mathrm dL=\ip{\Delta}{\mathrm dX\,W}_F+\ip{\Delta}{X\,\mathrm dW}_F
=\ip{\Delta W^{\trans}}{\mathrm dX}_F+\ip{X^{\trans}\Delta}{\mathrm dW}_F,
\]

两个梯度一次读出：\(\nabla_X L=\Delta W^{\trans}\)，\(\nabla_W L=X^{\trans}\Delta\)。整个过程没有出现一个下标，也没有用到任何需要背的矩阵求导公式；用到的只有微分规则、内积搬移和一次形状核对。换成注意力里的 \(QK^{\trans}\)、归一化层里的均值方差、低秩分解里的两个因子，步骤完全一样。这也是为什么框架的自定义算子文档总在讲``给定上游梯度，返回下游梯度''——那正是 \(\mathrm dL\) 被整理成内积之后括号里剩下的东西。

\subsection{和与平均只差一个因子，却改变有效步长}

目标函数写成样本损失之和还是平均，梯度会整体差一个 \(1/n\)。两份资料的公式若只差这个因子，多半各自对应了不同的目标定义，都成立。数学上没有对错，工程上要当心：同一个学习率下有效步长差了 \(n\) 倍，权重衰减相对于损失的比重也跟着变。换批大小、换成梯度累积、从复现代码里搬公式，这三处最容易埋雷。写下梯度之前先确认目标里有没有那个 \(1/n\)，比事后调学习率省事。

\subsection{变量有约束时要先确认扰动空间}

如果 \(X\) 被要求对称，那么可接受的 \(\mathrm dX\) 也必须对称，按无约束流程读出的 \(G\) 就不能直接当梯度用。原因可以算出来：设 \(K^{\trans}=-K\)、\(H^{\trans}=H\)，则

\[
\trace(KH)=\trace\bigl((KH)^{\trans}\bigr)=\trace(H^{\trans}K^{\trans})=-\trace(HK)=-\trace(KH),
\]

故 \(\trace(KH)=0\)。斜对称部分与任何对称扰动的内积都是零，对目标函数的变化没有贡献。留下负责的那一半：

\[
\operatorname{sym}(G)=\tfrac12\bigl(G+G^{\trans}\bigr),
\]

它自身对称，且对所有对称 \(H\) 有 \(\ip{G}{H}_F=\ip{\operatorname{sym}(G)}{H}_F\)。协方差矩阵、核矩阵、图的邻接权重上做优化时，漏掉这一步会让更新方向带上一个不属于问题的分量。要不要做这一步，取决于变量空间是怎么定义的，从公式本身看不出来。

到这里，单个函数的梯度已经可以机械地算出来。可实际的模型是几十层复合起来的，微分要一层层往回传，而传的时候有一个自由度还没用上：一串矩阵连乘，从哪一端开始乘是你说了算的。下一节就来算这笔账。
